\paragraph{IV 20.12: Généralisation de résultats antérieurs aux sections linéaires (cote 25, sites 66--74).}
Grothendieck extends the notation and the working constructions of the preceding hyperplane sections to linear sections of codimension $m$. Site 66 distinguishes the indices for dimension and codimension, introduces the Grassmannian parameter space and its incidence section, and relates the two transposed homomorphisms defining that section. At site 67 he explicitly criticizes the choice of a new letter $Y$ for the incidence construction and proposes $\underline{X}^{(m)}$; his displayed comparison of iterated exponents is accompanied by the words ``à peu de choses près''. The identification with a Grassmannian over $X$ leads into a detailed list of changes to the statements of sections 2--10 (sites 68--70), with several instructions addressed directly to Dieudonné. The added condition on dimensions, the repeated use of $F$-regularity, and the appeal to paragraph 19 belong to that list. A canceled initial objection at site 71 gives way to an extended discussion of the scope and editorial cost of the proposed generalization. The singular and supersingular loci $V^1$, $V''$, and $V^2$ are then related to the complements to paragraph 16 and to the filtration of a Grassmannian by dimensions of intersection with a fixed linear subspace (sites 71--73). His qualifications ``sauf erreur'' and ``sauf erreur de calcul'' accompany this discussion and its dimension expression; they remain visible in the text. Site 74 turns to the projective family $Q(C)$ of linear spaces through a fixed $C$, recalls 10.2 with a retouched following reference, and adds the previously omitted nonramification and Veronese-composition requirements to the proposed statement. The final handwritten continuation retains $Z$ in its last ``assez général'' phrase. Other notation requiring attention includes the unaccented $E$ followed by ``soit F'' in the incidence construction at site 66, the canceled material after the Grassmannian subscript $m$ on that page, and the baseline typed $v$ after $G_T$. These are readings of the scans, and are left in place; the apparatus records the written changes without supplying a mathematical revision.
